\begin{table}[H]
	\centering
	\caption{Spesifikasi Kebutuhan Non-Fungsional Model}
	\label{tab:kebutuhan_non_fungsional}
	\begin{tabular}{|c|l|p{9.5cm}|}
		\hline
		\textbf{Kode} & \textbf{Kebutuhan Non-Fungsional} & \textbf{Deskripsi} \\ \hline
		NF01 & \textit{Interpretability} & Model harus mampu menyediakan transparansi atas keputusan yang diambil. Berbeda dengan model sebelumnya yang bersifat \textit{black-box}, model yang diusulkan harus memvisualisasikan area fokus, seperti penggunaan \textit{heatmap} yang menjadi dasar model memprediksi. Hal ini bertujuan untuk menjawab keraguan pasien akan validitas diagnosis AI dan membantu profesional medis memverifikasi hasil tanpa harus mempercayai algoritma secara buta. \\ \hline
		NF02 & \textit{Reliability} & Model harus memiliki tingkat konsistensi yang tinggi dalam memberikan hasil prediksi. Model diharapkan dapat mempertahankan atau melampaui akurasi model referensi. \\ \hline
		NF03 & \textit{Complexity} & Model harus memiliki tingkat \textit{space complexity} yang rendah dan terukur pada tahap inferensi. Hal ini penting agar model dapat diimplementasikan secara optimal pada infrastruktur komputasi yang umum tersedia di fasilitas kesehatan, tanpa mengorbankan kecepatan maupun akurasi proses klasifikasi dan visualisasi. \\ \hline
	\end{tabular}
\end{table}